\documentclass[aspectratio=169,10pt]{beamer}

% Shared Beamer setup for Fall 2026 course slides.
% Clean Cal Poly-inspired theme: no custom class files or uncommon packages.
\mode<presentation>{
  \usetheme{default}
  \usefonttheme{professionalfonts}
}

\definecolor{CalPolyGreen}{HTML}{154734}
\definecolor{CalPolyGold}{HTML}{BD8B13}
\definecolor{CalPolySage}{HTML}{7FA28F}
\definecolor{CalPolyMint}{HTML}{DDF1E7}
\definecolor{CalPolySand}{HTML}{A29F86}
\definecolor{CalPolyBeige}{HTML}{EFEEDE}
\definecolor{DeckInk}{HTML}{1F2933}
\definecolor{DeckMuted}{HTML}{5F6B7A}
\definecolor{DeckSoft}{HTML}{F7F7F5}

\setbeamersize{text margin left=0.55in,text margin right=0.55in}
\setbeamercolor{normal text}{fg=DeckInk,bg=white}
\setbeamercolor{structure}{fg=CalPolyGreen}
\setbeamercolor{title}{fg=white}
\setbeamercolor{frametitle}{fg=CalPolyGreen,bg=white}
\setbeamercolor{block title}{bg=CalPolySage,fg=white}
\setbeamercolor{block body}{bg=CalPolyMint,fg=black}
\setbeamercolor{alerted text}{fg=CalPolyGold}

\setbeamerfont{title}{size=\Huge,series=\bfseries}
\setbeamerfont{frametitle}{size=\Large,series=\bfseries}
\setbeamerfont{block title}{size=\normalsize,series=\bfseries}

\setbeamertemplate{navigation symbols}{}
\setbeamertemplate{itemize item}{\color{CalPolyGold}$\blacktriangleright$}
\setbeamertemplate{itemize subitem}{\color{CalPolyGreen}$\bullet$}
\setbeamertemplate{footline}{
  \hfill{\color{DeckMuted}\scriptsize\insertframenumber/\inserttotalframenumber}
  \hspace{0.18in}\vspace{0.12in}
}
\setbeamertemplate{frametitle}{
  \vspace{0.18in}
  \hspace{0.02in}\insertframetitle\par
  \vspace{0.08in}
  \noindent\makebox[\linewidth][l]{%
    {\color{CalPolyGold}\rule{0.55in}{2pt}}%
    {\color{CalPolyGreen}\rule{\dimexpr\linewidth-0.55in\relax}{2pt}}%
  }
}

\newcommand{\CourseNumber}{Course}
\newcommand{\CourseTitle}{Course Title}
\newcommand{\LectureNumber}{0}
\newcommand{\LectureTitle}{Lecture Title}
\newcommand{\LectureDate}{Date}
\newcommand{\CourseUnit}{Unit}

\newcommand{\coursenumber}[1]{\renewcommand{\CourseNumber}{#1}}
\newcommand{\coursetitle}[1]{\renewcommand{\CourseTitle}{#1}}
\newcommand{\lecturenumber}[1]{\renewcommand{\LectureNumber}{#1}}
\newcommand{\lecturetitle}[1]{\renewcommand{\LectureTitle}{#1}}
\newcommand{\lecturedate}[1]{\renewcommand{\LectureDate}{#1}}
\newcommand{\courseunit}[1]{\renewcommand{\CourseUnit}{#1}}

\author{Kirk Duran}
\institute{California Polytechnic State University, San Luis Obispo}
\date{\LectureDate}

\newcommand{\makedecktitle}{
  \begingroup
  \setbeamercolor{background canvas}{bg=CalPolyGreen}
  \begin{frame}[plain]
    \vfill
    \begin{center}
      {\usebeamerfont{title}\color{white}\LectureTitle\par}
      \vspace{0.18in}
      {\Large\color{white}Lecture \LectureNumber\par}
    \end{center}
    \vfill
    \noindent{\color{white}\rule{\linewidth}{1.2pt}}\par
    \vspace{0.08in}
    \noindent\colorbox{CalPolyGold}{\makebox[0.24\linewidth][c]{\color{white}\Large\CourseNumber}}
    \hspace{0.12in}{\color{white}\Large Kirk Duran}
  \end{frame}
  \endgroup
}

\newenvironment{keyidea}{\begin{block}{Key idea}}{\end{block}}
\newenvironment{activity}{\begin{block}{In class}}{\end{block}}
\newenvironment{cpdefinition}{\begin{block}{Definition}}{\end{block}}
\newenvironment{exampleblockcp}{
  \begingroup
  \setbeamercolor{block title}{bg=CalPolySand,fg=white}
  \setbeamercolor{block body}{bg=CalPolyBeige,fg=black}
  \begin{block}{Example}
}{
  \end{block}
  \endgroup
}

\newcommand{\imageplaceholder}[2][0.3in]{%
  \noindent\makebox[\linewidth][c]{%
    \fcolorbox{CalPolySage}{CalPolyBeige}{%
      \parbox[c][#1][c]{0.86\linewidth}{%
        \centering
        {\color{DeckMuted}\scriptsize Image placeholder: }%
        {\color{DeckInk}\scriptsize\ttfamily\detokenize{#2}}%
      }%
    }%
  }%
}

\newcommand{\sectionframe}[1]{
  \begingroup
  \setbeamercolor{background canvas}{bg=CalPolyGreen}
  \begin{frame}[plain]
    \vfill
    {\color{white}\LARGE\bfseries #1\par}
    \vspace{0.18in}
    {\color{CalPolyGold}\rule{0.22\paperwidth}{3pt}}
    {\color{white}\rule{0.62\paperwidth}{3pt}}
    \vfill
  \end{frame}
  \endgroup
}

\newcommand{\placeholdernote}{
  \begin{block}{Placeholder}
    Replace these bullets with the final lecture material, examples, and in-class activities.
  \end{block}
}

\usepackage[T1]{fontenc}
\usepackage{lmodern}
\pdfstringdefDisableCommands{\def\translate#1{#1}}

\graphicspath{{../assets/lecture-29/}}

% Generated classroom visuals live in ../assets/lecture-29/.
\newcommand{\lectureimage}[2][1.75in]{%
  \includegraphics[width=\linewidth,height=#1,keepaspectratio]{#2}%
}

\setbeamersize{text margin left=0.42in,text margin right=0.42in}

\coursenumber{CS 1001}
\coursetitle{Introduction to Computing}
\lecturenumber{29}
\lecturetitle{Objects II --- Methods, Aliasing, and Collections}
\lecturedate{Fall 2026}
\courseunit{11. Objects and Compound Data}

\begin{document}

\makedecktitle

\begin{frame}{Today}
  \begin{itemize}
    \item Define \textbf{methods} as behavior associated with a class.
    \item Explain how \texttt{self} identifies the particular instance involved in a method call.
    \item Distinguish methods that return information from methods that mutate object state.
    \item Extend the reference and aliasing model to user-defined objects.
    \item Distinguish object identity from equality at a basic level.
    \item Store objects inside lists and dictionaries and iterate over those collections.
    \item Build objects that themselves contain collections or other objects.
  \end{itemize}

  \begin{keyidea}
    Methods operate on object state, while references determine \emph{which} object a method call or mutation affects.
  \end{keyidea}
\end{frame}

\begin{frame}{Where We Left Off}
  \begin{columns}[T,onlytextwidth]
    \begin{column}{0.48\textwidth}
      \begin{block}{Lecture 28}
        \begin{itemize}
          \item class and instance
          \item \texttt{\_\_init\_\_}
          \item attributes
          \item object state
          \item references to instances
        \end{itemize}
      \end{block}
    \end{column}
    \begin{column}{0.48\textwidth}
      \begin{block}{Lecture 29}
        \begin{itemize}
          \item behavior through methods
          \item deeper meaning of \texttt{self}
          \item mutation and aliasing
          \item identity and equality
          \item collections of objects
          \item nested compound structures
        \end{itemize}
      \end{block}
    \end{column}
  \end{columns}
\end{frame}

\sectionframe{Part 1: From Object State to Object Behavior}

\begin{frame}{Objects Can Represent Both State and Behavior}
  \begin{columns}[T,onlytextwidth]
    \begin{column}{0.53\textwidth}
      \begin{block}{State}
        A \texttt{Dog} instance might store:
        \begin{itemize}
          \item \texttt{name}
          \item \texttt{breed}
          \item \texttt{age}
        \end{itemize}
      \end{block}
      \begin{block}{Behavior}
        A \texttt{Dog} might also be able to:
        \begin{itemize}
          \item report a description
          \item compute a value from its state
          \item change its state
        \end{itemize}
      \end{block}
    \end{column}
    \begin{column}{0.40\textwidth}
      % Intended visual: one Dog instance with attribute fields on one side and method actions on the other.
      \lectureimage[2.22in]{object-state-and-behavior.png}
    \end{column}
  \end{columns}
\end{frame}

\begin{frame}{Definition: Method}
  \small
  \begin{cpdefinition}
    A \textbf{method} is a function defined as part of a class. A method can use the state of a particular instance and can optionally modify that state.
  \end{cpdefinition}

  \begin{block}{Example}
    \texttt{class Dog:}\\
    \hspace*{0.20in}\texttt{def \_\_init\_\_(self, name: str, breed: str, age: int):}\\
    \hspace*{0.40in}\texttt{self.name = name}\\
    \hspace*{0.40in}\texttt{self.breed = breed}\\
    \hspace*{0.40in}\texttt{self.age = age}\\
    \hspace*{0.20in}\texttt{def bark(self) -> str:}\\
    \hspace*{0.40in}\texttt{return self.name + " says bark!"}
  \end{block}

  \begin{itemize}
    \item Its first parameter is conventionally named \texttt{self}.
    \item It is called using dot notation on an instance.
  \end{itemize}
\end{frame}

\begin{frame}{Calling a Method Uses Dot Notation}
  \begin{block}{Class and call}
    \texttt{dog1 = Dog("Sofia", "Maremma", 1)}\\[0.08in]
    \texttt{message = dog1.bark()}
  \end{block}

  \begin{itemize}
    \item \texttt{dog1} refers to one particular \texttt{Dog} instance.
    \item \texttt{.bark} selects behavior defined by the \texttt{Dog} class.
    \item The method executes using \texttt{dog1} as the current instance.
    \item The returned string is bound to \texttt{message}.
  \end{itemize}

  \begin{keyidea}
    Dot notation connects a method definition to the instance on which that behavior should operate.
  \end{keyidea}
\end{frame}

\begin{frame}{What Does \texttt{self} Mean?}
  \begin{columns}[T,onlytextwidth]
    \begin{column}{0.55\textwidth}
      \begin{cpdefinition}
        Inside an instance method, \texttt{self} is the parameter that refers to the instance involved in the current method call.
      \end{cpdefinition}

      \begin{itemize}
        \item \texttt{self.name} means the \texttt{name} attribute of that instance.
        \item \texttt{self.age} means the \texttt{age} attribute of that instance.
        \item Different calls can bind \texttt{self} to different instances.
      \end{itemize}
    \end{column}
    \begin{column}{0.39\textwidth}
      % Intended visual: dog1.bark() arrow into a method frame where self points back to dog1.
      \lectureimage[2.22in]{self-method-binding.png}
    \end{column}
  \end{columns}
\end{frame}

\begin{frame}{Trace: \texttt{self} Is Bound for the Call}
  \begin{block}{Program}
    \texttt{dog1 = Dog("Sofia", "Maremma", 1)}\\
    \texttt{dog2 = Dog("Buddy", "Labrador", 3)}\\[0.08in]
    \texttt{a = dog1.bark()}\\
    \texttt{b = dog2.bark()}
  \end{block}

  \begin{center}
    \renewcommand{\arraystretch}{1.25}
    \begin{tabular}{l|l|l}
      \textbf{Call} & \textbf{\texttt{self} refers to} & \textbf{attribute read} \\\hline
      \texttt{dog1.bark()} & \texttt{dog1} instance & \texttt{dog1.name} \\
      \texttt{dog2.bark()} & \texttt{dog2} instance & \texttt{dog2.name}
    \end{tabular}
  \end{center}
\end{frame}

\begin{frame}{Methods Can Accept Additional Arguments}
  \begin{block}{Add this method inside the same Dog class}
    \texttt{def rename(self, new\_name: str) -> None:}\\
    \hspace*{0.20in}\texttt{self.name = new\_name}
  \end{block}

  Calling:
  \begin{block}{Call}
    \texttt{dog1.rename("Luna")}
  \end{block}

  conceptually binds:
  \begin{itemize}
    \item \texttt{self} $\rightarrow$ the instance referred to by \texttt{dog1}
    \item \texttt{new\_name} $\rightarrow$ \texttt{"Luna"}
  \end{itemize}
\end{frame}

\sectionframe{Part 2: Query Methods and Mutating Methods}

\begin{frame}{Some Methods Compute Information Without Changing State}
  \begin{block}{Add this query method inside Dog}
    \texttt{def human\_years(self) -> int:}\\
    \hspace*{0.20in}\texttt{return self.age * 7}
  \end{block}

  \begin{itemize}
    \item The method reads \texttt{self.age}.
    \item It computes and returns a new value.
    \item The object's attributes are unchanged.
  \end{itemize}

  \begin{keyidea}
    Reading object state does not imply mutating object state.
  \end{keyidea}
\end{frame}

\begin{frame}{Other Methods Intentionally Change State}
  \begin{block}{Add this mutating method inside Dog}
    \texttt{def birthday(self) -> None:}\\
    \hspace*{0.20in}\texttt{self.age = self.age + 1}
  \end{block}

  \begin{itemize}
    \item The right-hand side reads the current age.
    \item Assignment writes a new value into the instance's \texttt{age} attribute.
    \item The change remains after the method returns.
    \item The useful result is the changed object state, so no explicit value needs to be returned.
  \end{itemize}
\end{frame}

\begin{frame}{Query Versus Mutation Is an Interface Decision}
  \begin{center}
    \renewcommand{\arraystretch}{1.28}
    \begin{tabular}{l|l|l}
      \textbf{Method style} & \textbf{Primary effect} & \textbf{Typical return} \\\hline
      query & reads state & computed value \\
      mutation & changes state & often \texttt{None}
    \end{tabular}
  \end{center}

  \begin{block}{Examples}
    \texttt{dog1.human\_years()} $\rightarrow$ returns information\\
    \texttt{dog1.birthday()} $\rightarrow$ changes \texttt{dog1.age}
  \end{block}

  \begin{keyidea}
    A method's contract should make it clear whether callers should expect a returned value, a state change, or both.
  \end{keyidea}
\end{frame}

\begin{frame}{Trace the State Change}
  \begin{block}{Before}
    \texttt{dog1} $\rightarrow$ \texttt{Dog(name="Sofia", age=1)}
  \end{block}

  \begin{block}{Call}
    \texttt{dog1.birthday()}
  \end{block}

  \begin{block}{After}
    \texttt{dog1} $\rightarrow$ \texttt{Dog(name="Sofia", age=2)}
  \end{block}

  \begin{itemize}
    \item The variable \texttt{dog1} was not rebound.
    \item The same instance still exists.
    \item One attribute inside that instance changed.
  \end{itemize}
\end{frame}

\sectionframe{Part 3: References and Object Aliasing}

\begin{frame}{Assignment Still Does Not Automatically Copy}
  \begin{columns}[T,onlytextwidth]
    \begin{column}{.50\textwidth}
      \begin{block}{Program}
        \texttt{dog1 = Dog("Sofia", "Maremma", 1)}\\
        \texttt{dog2 = dog1}
      \end{block}
      Assignment copies the reference; it does not construct another Dog.
    \end{column}
    \begin{column}{.44\textwidth}
      \lectureimage[1.75in]{object-aliasing-two-names.png}
    \end{column}
  \end{columns}
  \begin{keyidea}
    Two names are aliases when they refer to the same object.
  \end{keyidea}
\end{frame}

\begin{frame}{Mutation Through Either Alias Reaches the Same Instance}
  \begin{block}{Program}
    \texttt{dog1 = Dog("Sofia", "Maremma", 1)}\\
    \texttt{dog2 = dog1}\\
    \texttt{dog2.birthday()}
  \end{block}

  \begin{block}{After the call}
    \texttt{dog1.age} $\rightarrow$ \texttt{2}\\
    \texttt{dog2.age} $\rightarrow$ \texttt{2}
  \end{block}

  \begin{itemize}
    \item The method call occurred through the name \texttt{dog2}.
    \item The changed state belongs to the shared instance, not to a variable name.
  \end{itemize}
\end{frame}

\begin{frame}{Constructing Again Produces an Independent Instance}
  \begin{block}{Program}
    \texttt{dog1 = Dog("Sofia", "Maremma", 1)}\\
    \texttt{dog2 = Dog("Sofia", "Maremma", 1)}
  \end{block}

  \begin{itemize}
    \item The attribute values happen to match.
    \item Two constructor calls produce two distinct instances.
    \item Mutating one does not mutate the other.
  \end{itemize}

  \begin{columns}[T,onlytextwidth]
    \begin{column}{0.48\textwidth}
      \centering
      \texttt{dog1} $\rightarrow$ Dog instance A
    \end{column}
    \begin{column}{0.48\textwidth}
      \centering
      \texttt{dog2} $\rightarrow$ Dog instance B
    \end{column}
  \end{columns}
\end{frame}

\begin{frame}{Identity Asks Whether Two References Reach the Same Object}
  \begin{block}{Example}
    \texttt{dog1 = Dog("Sofia", "Maremma", 1)}\\
    \texttt{dog2 = dog1}\\
    \texttt{dog3 = Dog("Sofia", "Maremma", 1)}
  \end{block}

  \begin{center}
    \renewcommand{\arraystretch}{1.25}
    \begin{tabular}{l|l}
      \textbf{Expression} & \textbf{Result} \\\hline
      \texttt{dog1 is dog2} & \texttt{True} \\
      \texttt{dog1 is dog3} & \texttt{False}
    \end{tabular}
  \end{center}

  \begin{keyidea}
    \texttt{is} asks about identity: are these two references to the exact same instance?
  \end{keyidea}
\end{frame}

\begin{frame}{Identity and Equality Answer Different Questions}
  \begin{columns}[T,onlytextwidth]
    \begin{column}{0.46\textwidth}
      \begin{block}{Identity}
        \texttt{a is b}

        Are \texttt{a} and \texttt{b} references to the same object?
      \end{block}
    \end{column}
    \begin{column}{0.46\textwidth}
      \begin{block}{Equality}
        \texttt{a == b}

        Should these two values be considered equivalent?
      \end{block}
    \end{column}
  \end{columns}

  \begin{itemize}
    \item Identity is built into the reference model.
    \item A class can define its own notion of value equality.
    \item Our \texttt{Dog} class has no custom equality rule: distinct instances compare unequal with \texttt{==}, even if their attributes match.
  \end{itemize}
\end{frame}

\begin{frame}{Passing an Object to a Function Can Also Create an Alias}
  \begin{block}{Function}
    \texttt{def celebrate(dog: Dog) -> None:}\\
    \hspace*{0.20in}\texttt{dog.birthday()}
  \end{block}

  \begin{block}{Call}
    \texttt{dog1 = Dog("Sofia", "Maremma", 1)}\\
    \texttt{celebrate(dog1)}
  \end{block}

  During the call, the parameter \texttt{dog} and the caller's name \texttt{dog1} refer to the same instance.

  \begin{keyidea}
    Parameter binding does not automatically copy mutable objects.
  \end{keyidea}
\end{frame}

\sectionframe{Part 4: Collections of Objects}

\begin{frame}{Define the Student Used in These Collections}
  \begin{block}{A class for the next examples}
    \texttt{class Student:}\\
    \hspace*{0.20in}\texttt{def \_\_init\_\_(self, name: str, age: int, gpa: float):}\\
    \hspace*{0.40in}\texttt{self.name = name}\\
    \hspace*{0.40in}\texttt{self.age = age}\\
    \hspace*{0.40in}\texttt{self.gpa = gpa}
  \end{block}
  \begin{itemize}
    \item This version stores name, age, and GPA; Lecture 28's design stored a major instead of age.
    \item Read constructor arguments using the definition in the current example.
  \end{itemize}
\end{frame}

\begin{frame}{A List Can Contain References to Objects}
  \begin{block}{Example}
    \texttt{students = [}\\
    \hspace*{0.20in}\texttt{Student("Alice", 20, 3.8),}\\
    \hspace*{0.20in}\texttt{Student("Bob", 21, 3.5),}\\
    \hspace*{0.20in}\texttt{Student("Charlie", 19, 3.9)}\\
    \texttt{]}
  \end{block}

  \begin{itemize}
    \item The list is one collection value.
    \item Each element refers to a \texttt{Student} instance.
    \item Each instance carries its own compound state.
  \end{itemize}
\end{frame}

\begin{frame}{Mental Model: Collection of References}
  \begin{columns}[T,onlytextwidth]
    \begin{column}{0.47\textwidth}
      \begin{block}{List positions}
        \texttt{students[0]}\\
        \texttt{students[1]}\\
        \texttt{students[2]}
      \end{block}
    \end{column}
    \begin{column}{0.46\textwidth}
      % Intended visual: list slots containing arrows to three separate Student instances with attributes.
      \lectureimage[1.75in]{list-of-object-references.png}
    \end{column}
  \end{columns}

  \begin{keyidea}
    The list organizes references; the objects contain the student-specific state.
  \end{keyidea}
\end{frame}

\begin{frame}{Iterating Over a List of Objects}
  \begin{block}{Example}
    \texttt{for student in students:}\\
    \hspace*{0.20in}\texttt{print(student.name)}
  \end{block}

  Each iteration:
  \begin{enumerate}
    \item binds \texttt{student} to the next list element;
    \item that element is a reference to a \texttt{Student} instance;
    \item dot notation reads an attribute of that instance.
  \end{enumerate}
\end{frame}

\begin{frame}{Sequence Patterns Still Apply to Objects}
  Suppose we want the names of students with GPA at least \texttt{3.5}.

  \begin{block}{Explicit loop}
    \texttt{names = []}\\
    \texttt{for student in students:}\\
    \hspace*{0.20in}\texttt{if student.gpa >= 3.5:}\\
    \hspace*{0.40in}\texttt{names.append(student.name)}
  \end{block}

  \begin{keyidea}
    Earlier loop patterns do not disappear. The values being processed are now compound objects rather than primitive values.
  \end{keyidea}
\end{frame}

\begin{frame}{A Dictionary Can Also Store Objects}
  \begin{block}{Example}
    \texttt{students\_by\_id = \{}\\
    \hspace*{0.20in}\texttt{101: Student("Alice", 20, 3.8),}\\
    \hspace*{0.20in}\texttt{205: Student("Bob", 21, 3.5)}\\
    \texttt{\}}
  \end{block}

  Then:
  \begin{block}{Access}
    \texttt{students\_by\_id[205].name}
  \end{block}

  \begin{itemize}
    \item The dictionary associates an ID key with an object reference.
    \item Dot notation continues from the selected object.
  \end{itemize}
\end{frame}

\begin{frame}{Aliasing Can Occur Through a Collection}
  \begin{block}{Program}
    \texttt{alice = Student("Alice", 20, 3.8)}\\
    \texttt{students = [alice]}\\
    \texttt{students[0].gpa = 4.0}
  \end{block}

  After the mutation:
  \begin{block}{Observation}
    \texttt{alice.gpa} $\rightarrow$ \texttt{4.0}
  \end{block}

  \begin{itemize}
    \item \texttt{alice} and \texttt{students[0]} both reach the same instance.
    \item The list does not contain an independent copy of the student.
  \end{itemize}
\end{frame}

\sectionframe{Part 5: Objects Containing Structured Data}

\begin{frame}{An Object Attribute Can Itself Be a Collection}
  \begin{block}{Example}
    \texttt{class Classroom:}\\
    \hspace*{0.20in}\texttt{def \_\_init\_\_(self, name: str):}\\
    \hspace*{0.40in}\texttt{self.name = name}\\
    \hspace*{0.40in}\texttt{self.students = []}
  \end{block}

  \begin{itemize}
    \item \texttt{name} refers to a string.
    \item \texttt{students} refers to a mutable list.
    \item The object's state can therefore include structured values.
  \end{itemize}
\end{frame}

\begin{frame}{Methods Can Mutate Collection Attributes}
  \begin{block}{Add this method inside Classroom}
    \texttt{def add\_student(self, student: Student) -> None:}\\
    \hspace*{0.20in}\texttt{self.students.append(student)}
  \end{block}

  \begin{block}{Use}
    \texttt{room = Classroom("180-0114")}\\
    \texttt{alice = Student("Alice", 20, 3.8)}\\
    \texttt{room.add\_student(alice)}
  \end{block}

  \begin{keyidea}
    A mutating method may change an attribute by mutating the collection that attribute refers to. Appending an object stores a reference to that same object; it does not create an automatic copy.
  \end{keyidea}
\end{frame}

\begin{frame}{An Object Can Also Contain Another Object}
  \begin{block}{Example}
    \texttt{class Address:}\\
    \hspace*{0.20in}\texttt{def \_\_init\_\_(self, street: str, city: str):}\\
    \hspace*{0.40in}\texttt{self.street = street}\\
    \hspace*{0.40in}\texttt{self.city = city}
  \end{block}

  \begin{block}{Person contains an Address reference}
    \texttt{class Person:}\\
    \hspace*{0.20in}\texttt{def \_\_init\_\_(self, name: str, address: Address):}\\
    \hspace*{0.40in}\texttt{self.name = name}\\
    \hspace*{0.40in}\texttt{self.address = address}
  \end{block}
\end{frame}

\begin{frame}{Nested Attribute Access Is Evaluated Left to Right}
  \begin{block}{Example}
    \texttt{addr = Address("123 Main St", "San Luis Obispo")}\\
    \texttt{person = Person("Alice", addr)}\\[0.08in]
    \texttt{person.address.city}
  \end{block}

  Evaluate in stages:
  \begin{enumerate}
    \item \texttt{person.address} produces the referenced \texttt{Address} object.
    \item \texttt{.city} then reads that object's \texttt{city} attribute.
  \end{enumerate}
\end{frame}

\begin{frame}{Nested Objects Extend the Same Reference Model}
  \begin{columns}[T,onlytextwidth]
    \begin{column}{0.48\textwidth}
      \begin{itemize}
        \item A variable can refer to an object.
        \item An attribute can refer to another object.
        \item A list element can refer to an object.
        \item A dictionary value can refer to an object.
      \end{itemize}
    \end{column}
    \begin{column}{0.45\textwidth}
      % Intended visual: Person -> Address, Classroom -> list -> Student references.
      \lectureimage[1.65in]{nested-object-reference-graph.png}
    \end{column}
  \end{columns}

  \begin{keyidea}
    The syntax changes, but the fundamental model remains names and containers holding references to values.
  \end{keyidea}
\end{frame}

\sectionframe{Part 6: Integrated Tracing}

\begin{frame}{Integrated Example}
  For this final trace, use a fresh, smaller Student definition (name and GPA only).
  \begin{block}{Classes}
    \texttt{class Student:}\\
    \hspace*{0.20in}\texttt{def \_\_init\_\_(self, name: str, gpa: float):}\\
    \hspace*{0.40in}\texttt{self.name = name}\\
    \hspace*{0.40in}\texttt{self.gpa = gpa}\\
    \hspace*{0.20in}\texttt{def set\_gpa(self, gpa: float) -> None:}\\
    \hspace*{0.40in}\texttt{self.gpa = gpa}
  \end{block}
\end{frame}

\begin{frame}{Predict Before Tracing}
  \begin{block}{Program}
    \texttt{a = Student("Alice", 3.5)}\\
    \texttt{b = a}\\
    \texttt{students = [a]}\\
    \texttt{students[0].set\_gpa(3.9)}
  \end{block}

  Predict:
  \begin{enumerate}
    \item How many \texttt{Student} instances exist?
    \item What does \texttt{a.gpa} evaluate to?
    \item What does \texttt{b.gpa} evaluate to?
    \item Is \texttt{a is students[0]} true or false?
  \end{enumerate}
\end{frame}

\begin{frame}{Reveal: One Shared Instance}
  Only one constructor call occurred, so there is one \texttt{Student} instance.

  \begin{block}{Three references reach this one Student}
    \centering
    \texttt{a}, \texttt{b}, \texttt{students[0]}
    $\longrightarrow$ \texttt{name="Alice", gpa=3.9}
  \end{block}

  Therefore:
  \begin{itemize}
    \item \texttt{a.gpa} $\rightarrow$ \texttt{3.9}
    \item \texttt{b.gpa} $\rightarrow$ \texttt{3.9}
    \item \texttt{a is students[0]} $\rightarrow$ \texttt{True}
  \end{itemize}
\end{frame}

\begin{frame}{A Useful Tracing Procedure for Objects}
  When object code becomes difficult, trace these questions in order:

  \begin{enumerate}
    \item Which instances have actually been constructed?
    \item Which names or container positions refer to each instance?
    \item Which instance is bound to \texttt{self} for the current method call?
    \item Which attributes are read?
    \item Which attributes or contained mutable values are changed?
    \item What value, if any, is returned?
  \end{enumerate}

  \begin{keyidea}
    Separate \textbf{control flow}, \textbf{references}, and \textbf{object state}; do not mentally collapse them into one step.
  \end{keyidea}
\end{frame}

\begin{frame}{Three Common Wrong Models}
  \begin{columns}[T,onlytextwidth]
    \begin{column}{0.31\textwidth}
      \begin{alertblock}{Names own objects}
        \texttt{b = a} does \textbf{not} construct a second object. It creates another reference to the same object.
      \end{alertblock}
    \end{column}
    \begin{column}{0.31\textwidth}
      \begin{alertblock}{Collections copy objects}
        Putting \texttt{alice} into a list does \textbf{not} automatically create an independent student.
      \end{alertblock}
    \end{column}
    \begin{column}{0.31\textwidth}
      \begin{alertblock}{\texttt{self} is global}
        \texttt{self} is a method parameter whose binding is established separately for each call.
      \end{alertblock}
    \end{column}
  \end{columns}
\end{frame}

\begin{frame}{What We Are Not Adding Yet}
  \begin{itemize}
    \item inheritance and class hierarchies
    \item private attributes or property decorators
    \item class variables and static methods
    \item custom operator behavior
    \item abstract classes or interfaces
    \item advanced object-oriented design patterns
  \end{itemize}

  \begin{keyidea}
    The goal is a precise mental model of instances, methods, references, and compound structures before adding larger object-oriented abstractions.
  \end{keyidea}
\end{frame}

\begin{frame}{Takeaways}
  \begin{enumerate}
    \item A method is behavior defined as part of a class.
    \item \texttt{self} refers to the particular instance involved in the current method call.
    \item Methods may query state, mutate state, or deliberately do both.
    \item Assignment and parameter binding can create aliases to the same object.
    \item \texttt{is} asks about object identity; equality is a separate concept.
    \item Collections can store references to objects, and objects can contain collections or other objects.
    \item The same reference model developed for lists now explains compound object structures.
  \end{enumerate}

  \begin{block}{Lab connection}
    Lab Programming 3 will ask you to develop functions that operate on objects using these same state, reference, and interface ideas.
  \end{block}
\end{frame}

\end{document}
